%%\documentclass[12pt,a4paper,oneside]{book}
\usepackage[ top=2.5cm,bottom=2.5cm,left=3cm,right=2cm ]{ geometry }

\usepackage{subcaption}
\usepackage{caption}

\usepackage[no-math]{fontspec}   %加這個就可以設定字體 
\usepackage{xeCJK}       %讓中英文字體分開設置 
\usepackage[UTF8,fontset=none,scheme=plain]{ctex}
\setCJKmainfont[AutoFakeBold=4, AutoFakeSlant=0.2]{TW-Kai}
\setCJKmonofont{Noto Sans Mono CJK TC}
\setmainfont{Times New Roman}
\usepackage{CJKnumb}

\usepackage{indentfirst}
\usepackage{array}
\usepackage{multirow}
\usepackage{amsthm}
\usepackage{amsmath}
\usepackage{amsfonts}
\usepackage{amssymb}
\usepackage{graphicx}
\usepackage{multirow}
\usepackage{setspace}
%\usepackage{cite}
\usepackage{balance}
\usepackage{textcomp}
\usepackage{color}
\usepackage[ampersand]{easylist}
\ListProperties(Hide=100, Hang=true, Progressive=3ex, Style*=$\bullet$ ,
Style2*=\tiny$\blacksquare$ ,Style3*=$\circ$ ,Style4*=-- )
\usepackage{changepage}
\usepackage{algorithm}
\usepackage{algpseudocode}
\usepackage{graphics}
\usepackage{epsfig}
\floatname{algorithm}{Algorithm}
\renewcommand{\algorithmicrequire}{\textbf{Input:}}
\renewcommand{\algorithmicensure}{\textbf{Output:}}
\renewcommand{\algorithmiccomment}{// }
\usepackage{verbatim}
\newtheorem{theorem}{Theorem}
\newtheorem{lemma}{Lemma}
\newtheorem{prop}{Proposition}
\newtheorem{defn}{Definition}

\usepackage{eqparbox}
\usepackage{textcomp}

\renewcommand\algorithmiccomment[1]{%
  \hfill\#\ \eqparbox{COMMENT}{#1}%
}

\usepackage{verbatim}

\usepackage{indentfirst}

\usepackage{chngcntr}
\counterwithout{figure}{chapter}
\counterwithout{table}{chapter}

\usepackage{xcolor}

\definecolor{codegreen}{rgb}{0,0.6,0}
\definecolor{codepurple}{rgb}{0.58,0,0.82}
\definecolor{backcolour}{rgb}{0.95,0.95,0.92}

% source code hightlighting
\usepackage{listings}
\lstset{
    captionpos=b,
    frame=single,
    keepspaces=true,                 
    numbers=left,                    
    numbersep=15pt,
    numberbychapter=false,
    lineskip=10pt,
    basicstyle=\ttfamily\footnotesize,
    tabsize=2,
    showstringspaces=false,
    language=C,
    keywordstyle=\color{magenta},
    backgroundcolor=\color{backcolour},   
    commentstyle=\color{codegreen},
    stringstyle=\color{codepurple},
    xleftmargin=.1\textwidth, xrightmargin=.1\textwidth,
    framexrightmargin=5pt, framexleftmargin=5pt,
    aboveskip=30pt, belowskip=30pt
}

% setting the page number to footer
\usepackage{fancyhdr}
\fancyhf{}
\cfoot{\thepage}
\pagestyle{fancy}
% no header and footer bar
\renewcommand{\headrulewidth}{0pt}
\renewcommand{\footrulewidth}{0pt}

% line height setting
\linespread{1.5}
\usepackage{setspace}

\usepackage{pdfpages}

% 浮水印套件 + 語法
\usepackage{eso-pic}
\newcommand\MyAtPageCenter[1]{\AtPageUpperLeft{%
\put(\LenToUnit{.42\paperwidth},\LenToUnit{-.5\paperheight}){#1}}%
}
% 浮水印套件 + 語法 end

% 調整標題格式
\usepackage{titlesec}
\usepackage{titletoc}

% 自訂語法
\usepackage{etoolbox}
\newtoggle{toc-use-cn}
\newtoggle{iamphd}
\newtoggle{EnableDynTitle}

\usepackage{adjustbox}
\usepackage{xstring}
\usepackage{anyfontsize}

% book table語法
\usepackage{booktabs}

\usepackage[flushleft]{threeparttable}
\usepackage{tabularx}
\usepackage{makecell}
